\paragraph{A relative fermion determinant for the doubleton member at fixed \(\gamma>0\) (conditional on a supplied trace and prescription; bounded result).}
Keep the setting of the preceding paragraph. In the one-loop formula of~\cite{Steinacker2023} (its Eqs.~(4.15)--(4.17)) the fermions enter through \(-\tfrac12\operatorname{Tr}\log(\Box-\mathrm i\epsilon-\Sigma^{(\psi)}_{ab}[\mathcal F^{ab},\cdot\,])\) with \(\Box=\eta^{ab}[\bar T_a,[\bar T_b,\cdot\,]]\), \(\mathcal F^{ab}=\mathrm i[\bar T^a,\bar T^b]\) and \(\Sigma^{(\psi)}_{ab}=\tfrac{\mathrm i}4[\Gamma_a,\Gamma_b]\). Since \((\Gamma^a[\bar T_a,\cdot\,])^2=\Box-\Sigma^{(\psi)}_{ab}[\mathcal F^{ab},\cdot\,]\), this is the logarithm of the square \(D^2\) of the fermion operator. The \(\mathrm i\epsilon\) of that reference comes from a Gaussian factor on the bosons (its Eq.~(4.15)), so in the fermion logarithm it is a prescription. It is supplied here, as are the branch and the trace. Define
\[
\log\operatorname{Pf}_{\mathrm{rel}}=\tfrac14\,\tau\bigl[\operatorname{tr}_{16}\log(D^2-\mathrm i\epsilon)-16\log(\Box-\mathrm i\epsilon)\bigr],
\]
with the principal branch and the spinor trace taken first. Write the Hilbert--Schmidt operators as \(\int\pi_{k,p}\otimes\mathcal M_{k,p}\,\mathrm dm(k,p)\) over the unitary principal series of \(\mathrm{Spin}(3,1)\), as in the preceding paragraph; \(k\) is an integer, since \(-1\in\mathrm{Spin}(3,1)\) is central in the group acting on the irreducible \(\mathcal H_n\), acts there by a scalar and hence trivially in the adjoint action. Then \(\tau=\int\mathrm dm(k,p)\,\omega_{k,p}\), with \(\omega_{k,p}\) a normal semifinite weight on the von Neumann algebra generated by \(C_5\), which acts on \(\mathcal M_{k,p}\), that is, a measure on the spectrum of \(C_5\) there; only this restriction enters below. The trace of the operators on \(\mathcal M_{k,p}\) gives a spectral projection of \(C_5\) its rank, which is finite on an interval only if \(C_5\) has there finitely many eigenvalues of finite multiplicity. The weights \(10-8-2=0\) of that formula make its split into a relative bosonic and a relative fermionic term exact. Both \(D^2\) and \(\Box\) are proportional to \(c\), so for \(\epsilon\to0\) the fermionic term does not depend on \(\gamma>0\).

\emph{Fibre.} On one chirality \(c^{-1}\mathcal D_-\mathcal D_+=A-2X\), with \(A=C_5-C_4\) commuting with \(X\). On \(\pi_{k,p}\otimes\mathbb C^2\) one has \(X^2+X=\tfrac14(z^2-1)\), with \(z=k+\mathrm ip\) and its conjugate for the other chirality. So \(X\) has the roots \(\tfrac12(z-1)\) and \(-\tfrac12(z+1)\), and the spinor trace of a function \(f\) of \(\mathcal D_-\mathcal D_+\) is \(w_1f(\lambda_1)+w_2f(\lambda_2)\), with
\[
\lambda_{1,2}=c\bigl(C_5+2-k(k\pm1)+p^2\mp\mathrm ip\bigr),\qquad w_{1,2}=1\pm z^{-1}
\]
(exact). On finite-dimensional representations of the real form, with \(z=2j_{\mathrm L}+1\), the weights are the multiplicities per orbital dimension (checked on the full spectrum for \(n=1,2\)). The operator therefore depends on \((k,p)\) and on \(C_5\) only, and the distribution of \(C_5\) on \(\mathcal M_{k,p}\) enters only through \(\omega_{k,p}\).

\emph{Zeros and integrability.} One has \(\operatorname{Im}\lambda_{1,2}=\mp cp\). So the \(\lambda_i\) vanish only at \(p=0\) and \(C_5=k(k\pm1)-2\), which is the kernel condition of the preceding paragraph, and \(|\lambda_i-\mathrm i\epsilon|\ge|cp\pm\epsilon|\). The Plancherel density is bounded on bounded sets and vanishes to second order at \(z=0\), so \(w_i\,\mathrm dm\) is bounded there. Suppose \(\omega_{k,p}\) gives finite weight to the spectral projections of \(C_5\) for bounded intervals, locally uniformly in \((k,p)\). Then \(\log|\lambda_i-\mathrm i\epsilon|\) and its square are integrable on every bounded set of labels, uniformly in \(\epsilon\), and the arguments are bounded. By uniform integrability (Vitali's theorem) the fermionic term then has its limit \(\epsilon\to0\) there, and its zeros produce no divergence on such a set. The reference term has its limit there by dominated convergence if, in addition, \(\omega_{k,p}\) has a density in \(C_5\) bounded locally uniformly, since \(|\Box-\mathrm i\epsilon|\ge|\Box|\); under the first condition alone it is not controlled, since \(\Box=c(C_5-k^2+p^2+1)\) vanishes on a set that the weights may charge.

On the spectral subspace \(C_5\le-\tfrac94\) the bound is uniform, \(|\lambda_i|\ge c/4\), since there \(C_5+2-k(k\pm1)\le-\tfrac14\) for integer \(k\) (exact). This subspace contains the unitary spherical principal series of \(\mathrm{Spin}(4,1)\). There \(C_5=-\tfrac94-\nu^2\), and for \(k=0\) one has \(\lambda_1=c[(p-\tfrac{\mathrm i}2)^2-\nu^2]\). For \(-\tfrac94<C_5<-2\) the bound is \(c|C_5+2|\). On the principal series \((s,\nu)\), with \(s\) an integer for the same reason as \(k\), where \(C_5=s(s+1)-\tfrac94-\nu^2\), zeros therefore need \(s\ge1\).

\emph{Compact members.} For the compact members the same formula, with \(z\) real, puts the zero eigenvalue on the top component \((k_1,k_2)=(m_1,m_2)\) of the preceding paragraph but one. Its multiplicity is \((m_1+m_2+2)(m_1-m_2+1)\) per chiral block, which gives the kernel dimension \(K=\tfrac43n(n+4)(n^2+4n+7)\) on the sixteen components. The operator \(\Box\) has no kernel on the traceless matrices, so \(|\operatorname{Pf}_{\mathrm{rel}}|\) is proportional to \(\epsilon^{K/4}\) as \(\epsilon\to0\) and tends to zero (exponents \(20\) and \(76\) for \(n=1,2\), checked). In the labels above these zeros lie at imaginary \(p\) and \(\nu\), outside the unitary principal series.

\emph{Not established.} The set \(|\lambda|<\delta\) has finite \(\tau\)-measure only if the weights \(\omega_{k,p}\) of \(C_5\) near \(k(k\pm1)-2\) at small \(p\), multiplied by the Plancherel density, are summable over \(k\). That is a condition on the higher-spin content of the operators, and it is not checked here.

For large labels the first order of \(\sum_iw_i\log(\lambda_i/cA)\) in \(A^{-1}\) cancels. The second order is \((1-z^2)/A^2\), and summed over the chiralities it is \(-2C_4/A^2\). So the fermionic term alone shows no cancellation. At fixed \(C_5\) and \(k\), its density times the Plancherel density tends to a constant as \(p\to\infty\), by the second order above, so the trace at fixed \(C_5\) diverges linearly in a cut-off on \(p\) wherever \(\omega_{k,p}\) gives that value of \(C_5\) a weight bounded below at large \(p\). A trace per unit volume at a given time is not constructed. The bosons and ghosts are not included; in~\cite{Steinacker2023} they cancel the low orders against this term, and that cancellation fails at order \(\gamma\) (preceding paragraph but one).

Because \(\operatorname{Im}w_i\neq0\), the modulus of \(\operatorname{Pf}_{\mathrm{rel}}\) depends on the branch of the fermion logarithm. Its phase depends on the \(\mathrm i\epsilon\) of the reference \(\Box\).

So neither a finite non-zero density nor an infrared divergence is established. What is shown is that, under the conditions above and on every bounded set of labels, the zeros of the doubleton member, unlike the zero modes of the compact members, give a finite contribution to the fermionic term; whether these contributions are summable over \(k\) is the open condition stated above.

These statements are for one loop, \(N\to\infty\) first at fixed \(\gamma>0\) and then \(\epsilon\to0\), and the class of the preceding paragraph. For \(n>0\) they are conditional on the identification used there. Beyond one loop the compact zero modes are absorbed by the Yukawa vertices, so a vanishing one-loop factor does not remove those members. The family, \(h_0\), the mass, the fermion content, the trace and the prescription are supplied, so Problem~E stays open with Assumption~\ref{ass:lorentzian} retained.
